\chapter{Trabajo futuro}
\label{cap:trabajo-futuro}

% TODO: redactar trabajo futuro a partir de las conclusiones finales
% (capitulo 8) y de las recomendaciones consolidadas en la revision
% cruzada (capitulo 6). No anticipar aqui recomendaciones que todavia no
% fueron validadas en esos capitulos.

\pendienteIntegracion{Trabajo futuro, derivado de las conclusiones (cap\'itulo~\ref{cap:conclusiones}) y de las recomendaciones consolidadas de la revisi\'on cruzada (cap\'itulo~\ref{cap:revision-cruzada}).}

\section{Alcance previsto}
\label{sec:trabajo-futuro-alcance}

Esta secci\'on listar\'a las recomendaciones de mejora priorizadas que
surjan de la revisi\'on cruzada y de las conclusiones finales, diferenciando
expl\'icitamente entre:

\begin{itemize}[nosep]
  \item Mejoras t\'ecnicas concretas sobre BIOPET (por ejemplo, las
    limitaciones ya documentadas en el propio \texttt{README.md}: cach\'e
    de rendimiento, SEO en Lighthouse, cobertura de rama).
  \item Recomendaciones de proceso o documentaci\'on derivadas de la
    revisi\'on cruzada.
  \item L\'ineas de investigaci\'on adicional sugeridas por la comparaci\'on
    SOAP vs REST o por el an\'alisis de tendencias web, si aplican al
    contexto real de BIOPET.
\end{itemize}

No se redacta contenido definitivo en esta versi\'on del informe porque
depende de los cap\'itulos de conclusiones y revisi\'on cruzada, todav\'ia
pendientes de integraci\'on.
